%% ma: L10-B12-0023
%% lop: 10
%% chuong: 5
%% bai: 12
%% dang: 10.12.8
%% loai: TN4
%% muc: VD
%% dap_an: "C"
%% nguon: "chuyen-de-cac-so-dac-trung-cua-mau-so-lieu-khong-ghep-nhom-toan-10-chuong-trinh-moi.pdf (GV. Phan Nhật Linh) (P01), Bài 01 Phần I Câu 3"
%% kien_thuc: "Lớp 10 Bài 12 (số gần đúng, sai số tuyệt đối, độ chính xác, sai số tương đối, quy tròn số)"
%% trang_thai: chinh-thuc
%% ghi_chu: "Đề gốc hỏi 'cách viết chuẩn' (dạng chuẩn, chữ số chắc: không có trong SGK KNTT) với đáp án 26,7 sai; đã đổi thành hỏi số quy tròn theo SGK KNTT, đáp án 26,8; phương án D gốc 'Đáp án khác' đổi thành 26,75. Tự duyệt theo yêu cầu giáo viên 22/10/2026."
\begin{ex}
Đường kính của một đồng hồ cát là $8{,}52$ m với độ chính xác đến $1$ cm. Dùng giá trị gần đúng của $\pi$ là $3{,}14$ để tính chu vi $C=\pi d$ của đồng hồ. Số quy tròn của chu vi (theo độ chính xác tìm được của $C$) là
\choice{$26{,}6$}{$26{,}7$}{\True $26{,}8$}{$26{,}75$}
\loigiai{$8{,}51\le d\le8{,}53$ nên $26{,}7214\le C\le26{,}7842$, tức là $C=26{,}7528\pm0{,}0314$.

Chữ số khác $0$ đầu tiên của $0{,}0314$ ở hàng phần trăm, nên quy tròn $C$ đến hàng phần mười: $26{,}8$. Chọn C.}
\end{ex}
